\documentclass[11pt,a4paper]{article}
\usepackage[utf8]{inputenc}
\usepackage[T1]{fontenc}
\usepackage{lmodern}
\usepackage[french]{babel}
\usepackage{amsmath,amssymb,mathtools}
\usepackage{icomma}
\usepackage[version=4]{mhchem}
\usepackage{siunitx}
\sisetup{output-decimal-marker={,},per-mode=symbol}
\usepackage{xcolor}
\usepackage{colortbl}
\usepackage{tikz}
\usepackage[most]{tcolorbox}
\usepackage{enumitem}
\usepackage{array}
\usepackage{booktabs}
\usepackage{fancyhdr}
\usepackage[hidelinks]{hyperref}
\usetikzlibrary{arrows.meta,positioning,calc}
\usepackage[a4paper,margin=2.1cm,headheight=26pt,headsep=10pt]{geometry}

\definecolor{primary}{RGB}{146,95,161}
\definecolor{primarydark}{RGB}{92,55,108}
\definecolor{acidecol}{RGB}{205,60,45}
\definecolor{basecol}{RGB}{40,110,200}
\definecolor{methcol}{RGB}{90,90,100}

\tcbset{boxbase/.style={enhanced,breakable,boxrule=0.9pt,arc=2.5pt,left=3mm,right=3mm,top=2.3mm,bottom=2.3mm,fonttitle=\bfseries,coltitle=white,toptitle=0.6mm,bottomtitle=0.6mm}}
\newtcolorbox[auto counter]{exobox}[1][]{boxbase,colback=primary!4,colframe=primary,title={\textsc{Exercice}~\thetcbcounter\ifx\relax#1\relax\relax\else\textmd{ --- \itshape #1}\fi}}
\newtcolorbox{consigne}{boxbase,colback=methcol!7,colframe=methcol,sharp corners}

\pagestyle{fancy}
\fancyhf{}
\fancyhead[L]{\small\textcolor{primarydark}{\textbf{Fiche 10} — Thème 2 : pH, pOH, $K_e$}}
\fancyhead[R]{\small\textcolor{primarydark}{\textsc{Exercices — Difficile}}}
\fancyfoot[C]{\small\thepage}
\renewcommand{\headrulewidth}{0.8pt}
\renewcommand{\headrule}{\hbox to\headwidth{\color{primary}\leaders\hrule height \headrulewidth\hfill}}

\setlength{\parindent}{0pt}
\setlength{\parskip}{3pt}

\begin{document}

\begin{center}
{\LARGE\bfseries\color{primarydark} Thème 2 --- Niveau Difficile}\\[2pt]
{\color{primary}\itshape Température et pH neutre, comparaisons physiologiques, raisonnement sur $K_e$}
\end{center}

\begin{consigne}
\textbf{Objectif du niveau.} Combiner plusieurs conversions et raisonner sur les \emph{inégalités} entre
$[\ce{H3O+}]$ et $[\ce{OH-}]$. Chaque exercice se construit en sous-questions. \emph{On donne
$\log 2\approx0{,}30$ ; $\log 3\approx0{,}48$.}
\end{consigne}

\begin{exobox}[le pH neutre se déplace avec la température]
On chauffe de l'eau pure jusqu'à une température $T$ où le produit ionique vaut $K_e = 1{,}0\times10^{-13}$.
\begin{enumerate}[label=\alph*),leftmargin=1.6em,itemsep=2pt]
  \item Rappeler pourquoi $K_e$ augmente lorsque la température augmente (nature énergétique de
        l'autoprotolyse).
  \item Calculer $[\ce{H3O+}]$ dans cette eau pure (milieu neutre : $[\ce{H3O+}]=[\ce{OH-}]$), puis le
        $\mathrm{pH}$ neutre à cette température.
  \item À cette même température $T$, on mesure une solution à $\mathrm{pH}=7$. Est-elle acide, neutre ou
        basique ? Et une solution à $\mathrm{pH}=6{,}5$ ?
  \item En une phrase, expliquer pourquoi l'affirmation « une solution basique a toujours un pH supérieur à
        7 » est fausse en toute rigueur.
\end{enumerate}
\end{exobox}

\begin{exobox}[du suc gastrique au sang]
On compare deux fluides du corps humain, à \SI{25}{\celsius} pour simplifier : le suc gastrique
($\mathrm{pH}=1{,}5$) et le sang ($\mathrm{pH}=7{,}4$).
\begin{enumerate}[label=\alph*),leftmargin=1.6em,itemsep=2pt]
  \item Calculer $[\ce{H3O+}]$ pour chacun des deux fluides.
  \item Par quel facteur le suc gastrique est-il plus concentré en \ce{H3O+} que le sang ? (raisonner d'abord
        avec $10^{\Delta\mathrm{pH}}$, puis donner l'ordre de grandeur).
  \item Calculer $[\ce{OH-}]$ dans le sang à l'aide de $K_e$. Vérifier la cohérence avec $\mathrm{pOH}$.
\end{enumerate}
\end{exobox}

\begin{exobox}[{démontrer le critère « acide $\Rightarrow [\ce{H3O+}]>[\ce{OH-}]$ »}]
On se place à \SI{25}{\celsius} et on cherche à \emph{prouver}, sans se contenter de l'apprendre, qu'une
solution acide vérifie $[\ce{H3O+}]>[\ce{OH-}]$.
\begin{enumerate}[label=\alph*),leftmargin=1.6em,itemsep=2pt]
  \item Traduire « la solution est acide » par une inégalité sur le $\mathrm{pH}$, puis, en appliquant
        $[\ce{H3O+}]=10^{-\mathrm{pH}}$, en une inégalité sur $[\ce{H3O+}]$.
  \item À partir de $[\ce{H3O+}] > 10^{-7}$ et de $[\ce{OH-}] = K_e/[\ce{H3O+}]$, montrer que
        $[\ce{OH-}] < 10^{-7}$.
  \item Conclure : comparer $[\ce{H3O+}]$ et $[\ce{OH-}]$. Ce raisonnement resterait-il valable à une
        température où $K_e \neq 10^{-14}$ ? (répondre à propos de la valeur « 7 »).
\end{enumerate}
\end{exobox}

\end{document}
